\documentclass[10pt,a4paper]{amsart}
\usepackage[margin=19mm]{geometry}
\usepackage{amsmath,amssymb,mathtools}
\usepackage[scr=rsfs,cal=euler]{mathalfa}
\usepackage{xfrac}
\usepackage[unicode,hidelinks,bookmarks=false]{hyperref}
\newcommand{\ie}{i.e.\ }
\newcommand{\cf}{cf.\ }
\newcommand{\resp}{respectively\ }
\newcommand{\Subsection}{\subsection}
\providecommand{\nobreakdash}{\nobreak\hbox{-}}
\setlength{\emergencystretch}{2em}
\multlinegap0pt
\usepackage{bm}
\usepackage{amssymb}
\usepackage{xcolor}
\usepackage{enumerate}
\newtheorem{thm}{Theorem}
\newcommand{\e}{\varepsilon}
\newcommand{\es}{\emptyset}
\title{On best constrained and best simultaneous approximation in Banach spaces}
\author{Syamantak Das}
\date{Source: arXiv:2505.10851v3 (CC BY 4.0)}
\begin{document}
\maketitle

\noindent\emph{Source and licence.} On best constrained and best simultaneous approximation in Banach spaces. Version arXiv:2505.10851v3 (CC BY 4.0), distributed by arXiv under the Creative Commons Attribution 4.0 International licence, http://creativecommons.org/licenses/by/4.0/, as recorded in the arXiv licence metadata for this exact version and shown on the abstract page https://arxiv.org/abs/2505.10851v3. Full licence notice: \texttt{licenses/arxiv-2505.10851-CC-BY-4.0.txt}.\par\smallskip

\noindent\textit{Editorial scope.} Counted unit: the article's thm T3 (source0.tex lines 371--373). The article's complete original proof of it is delivered with it (source0.tex lines 374--382, one \texttt{proof} environment of 9 lines). No mathematics is written, repaired or re-derived: the statement and the proof are the article's own lines, in the article's own notation, and no retained line is substituted or re-flowed.  No further source-local context is needed: every cross-reference of every retained line resolves inside the two delivered spans.  Nothing was omitted from any retained span, so the package carries no omission record.  No retained line contains a citation, so the wrapper carries no bibliography and no bibliography entry.  No retained line contains a figure environment, an image-inclusion command or drawing code, so no image is delivered and the package declares no asset dependency.  Editorial formatting only, in addition: an AMS \texttt{amsart} preamble that restates the article's own declarations and macros used by the retained lines, namely the environments \texttt{thm} on separate counters, together with the standard packages those lines need.  The retained environments print in this document as Theorem 1 in order of appearance, whereas the article prints the same items with the numbering of its own full document.  The complete native source of the article as distributed in the arXiv e-print for this exact version is bundled unchanged as \texttt{sources/178204/source0.tex}.
\begin{thm}\label{T3}
Let $Y,Z_1$ be subspaces of $X$ and $Z_2=Z_1\cap Y$. Suppose $Z_2$ is locally constrained w.r.t. $Z_1$ via $Y$. Then if $Z_1$ is a (semi) $M$-ideal in $X$, then $Z_2$ is a (semi) $M$-ideal in $Y$.
\end{thm}

\begin{proof}
We prove the result for $M$-ideals. The proof for semi $M$-ideals is similar.

Let $(B_Y(a_i,r_i))_{i=1}^3$ be closed balls in $Y$ such that $\cap_{i=1}^3B_Y(a_i,r_i)\neq\es$ and $B_{Y}(a_i,r_i)\cap {Z_2}\neq\es$ for all $i=1,2,3$. As $Z_1$ is an $M$-ideal in $X$, for $\e>0$, there exists $z_0\in Z_1$ such that $z_0\in\cap_{i=1}^3B_X(a_i,r_i+\e)\cap Z_1$. Let $P,Q$ be projections for $z_0$ as defined in the definition of locally constrained spaces. Then, $Pz_0\in Z_2$ and for all $i=1,2,3$, we have
\[
\|Pz_0-a_i\|=\|Qz_0-a_i\|=\|Q(z_0-a_i)\|\leq\|z_0-a_i\|\leq r_i+\e.
\]
This concludes the proof.
\end{proof}

\end{document}
